% ----------------------------------------------------------------------
% ----------------------------------------------------------------------

\documentclass{hbarticle}
\usepackage{graphicx}
\usepackage{theorem}
\usepackage{amsfonts}
\usepackage{amssymb}
\usepackage{subfiles}
\usepackage{xcolor}
\usepackage{fancyhdr}

\DeclareRobustCommand{\mdelta}{\boldsymbol\delta}

%%%%%%%%%%%%%% Layoutgroesse%%%%%
\topmargin=-1cm
\textheight=24cm
\oddsidemargin=-0.2truecm
\evensidemargin=-0.2truecm
\textwidth=16cm%\textwidth=17cm
\renewcommand{\baselinestretch}{1.1}\normalsize
\overfullrule=0pt
%%%%%%%%%%%%%%%%%%%%%%%%%%%%%

\newcommand{\symbolfootnote}{\fnsymbol{footnote}}

\title{Discrete Time Risk Neutral Dynamics}
\author{Hans Buehler}

\pagestyle{fancy}
\lhead{Discrete Time Risk Neutral Dynamics}
\rhead{Buehler, Aug 2026}

\makeindex

\begin{document}


    \maketitle
    
    \begin{abstract}
        Finding a risk-neutral measure for discrete time
        markets.
    \end{abstract}

    \tableofcontents

    \section{Theory}
    
    We revisit the concept of risk-neutrality in markets with many
    more instruments that market factors relevant for modern generative
    machine learning applications in finance.
    
    To this end, our basic assumption is that our market dynamics
    is given over time steps~$t_j$ by a hidden state process~$h_j\in \calH \subseteq \R^d$
    which follows
    \eq{h}
    	h_{j+1} = G(h_j; Y_j)   
    \eqend
    where the generator $G$ of the state~$h_j$ 
    and some normal noise $Y_j\in \R^m$ defined
    under the real market (``statistical'') measure~$\P$. Essentially
    the latter means that this is our model learned to historic market data.


       
    We here explicitly assume
    that $\Delta_j := t_{j+1} - t_j$
    is not in the limit zero. Typical values are
    a day or several days. We will use '$dt$' to refer to infinisemal
    time increments.
   
   	\begin{example}[AR(1)]
    As an example, consider a classic
    diffusion AR(1) process like the one used in~\cite{DYSANOS}
    given in continuous time formulation as
    \eqx
	    	dh_t = ( \theta - K h_t )\, dt + \Sigma Y_t \sqrt{dt} \ .  
    \eqend
    where $\theta\in \R^d$, $K\in \R^{d,d}$ must have eigenvalues
    with non-negative 
    real values, and $\Sigma \in \R^{d,m}$ for $m\leq d$.
    
    In discrete-time increments of this process over~$\Delta t$ 
    are distributed exactly as
    \eqx
    	h_{j+1} = e^{-K\,\Delta_j} h_j + \int_0^{\Delta_j} 
    	e^{-K\,s} \bar h\,ds
    		+
    			S_\Delta \tilde Y_j  
    \eqend
    with
    \eqx		
  	   S_\Delta S _\Delta  = \int_0^{\Delta _j} ( e^{-K\,dt} \Sigma )(\Sigma' e^{-K'\,dt}) dt \ .
    \eqend
    Here $\tilde Y_t\in \R^{\mathrm{rank} S_\Delta}$
    with $\mathrm{rank} S_\Delta \in \{m,\ldots,d\}$.
    
    In the special case where all eigenvalues of $K$
    have positive real values, $\mathrm{rank} S_\Delta=d$
    and $K$ is invertible, which then leads to the 
    well-known \emph{flow} 
    \eqx
    	h_{j+1} = e^{-K\,\Delta_j} h_j + K^{-1}
    	\big( I - e^{-K\,\Delta_j} \big) \theta 
    	+
    	e^{-K\,s} \bar h\,ds
    		+
    			S_\Delta \tilde Y_j \ .
    \eqend
    This is a simple and transparent example of a process with in infinitesimal time
    $dt$ has a stochastic dimension of only $m\ll d$, but once discrete time steps
    are taken, its stochastic dimension increases.
    \end{example}
    
    We now assume that at each time step we are able to trade
    the stock price $S_j=\exp h^0_j$ and a large number of $L$ 
    of options, which are identified by their abstract
    contract \emph{terms}~$\xi^\ell\in \Xi$ which are shared accross all~$j$.  In~\cite{DYSANOS} for example
    this was simply~$\xi^\ell = (\tau^\ell, x^\ell)$
    where $\tau_j$ is the time-to-expiry and $x^\ell$ a log-strike price,
    with by convention $\xi^0 \equiv (0,0)$ for spot.
    All options were calls.
    We will assume in this article that all terms remain the same across $j$.
    We will also assume that all terms contain a time-to-expiry~$\tau^\ell$
    which $\infty$ for perpetual assets such as equities and FX.\footnote{
        In our agents we encode such expiries as zero.
    }
    
    The terms translate into a concrete \emph{instance} of a contract
    by applying them to the current environment represented by $h_j$.
    Formally each $\xi^\ell$
    is a map from $\calH$ into the set of $L^1$ integrable random variables. That yields the \emph{instances}
    $$
         \xi^\ell_j := \xi^\ell(h_j) \ .
    $$
    In above
    example the \emph{instance} for $\xi^\ell = (\tau^\ell, x^\ell)$
    is given by
    $$
        \xi^\ell_j := \left( \frac{ S_{t_j+\tau^\ell}}{ S_{t_j} } - x^\ell \right)^+ \ .
    $$
    We may think of this as the ``issuance'' of the contract; 
    \begin{assumption}
        Fix $h_0\in \calH $.  Then we can compute a finite mid-price
        $$
            H^\ell(h_0):= H\big(\, \xi^\ell(h_0); h_0, 0\, \big)
        $$
        and half-spread
        $$
            \gamma^\ell(h_0):= \gamma\big(\, \xi^\ell(h_0); h_0, 0\, \big)
        $$
        for all contracts $\xi^\ell\in\Xi$.      

        Let, moreover, $\tau > 0$ be a time step and $\tau^\ell$ the time-to-expiry for the $\ell$th option.
        We also assume that we may compute
        the mid-price and half-spread \emph{of the same instance} $\xi^\ell(h_0)$ as
        $$
            H^\ell_\tau(h_0):= H\left(\, \xi^\ell(h_0); h_{\min\{ \tau, \tau^\ell \}},\tau\, \right) \ \ \mbox{and}\ \
            \gamma^\ell_\tau(h_0):= \gamma\!\left(\, \xi^\ell(h_0); h_{\min\{ \tau, \tau^\ell \}},\tau \, \right)
        $$
        This condition simply means that either we can compute the expiry value of the option if it
        expired before~$\tau$, or that we can can compute its mid-price and half-spread from
        the current market, represented in its entirety by~$h_\tau$.
    \end{assumption}

    \begin{remark}
        We note that this
        formulation assumes that all our contracts require only one expiry value; it is easily extended
        to pathwise payoffs; it just clutters notation. The main point is that all future values of
        any tradable instrument depend
        only on future values of the Markov process~$h$.
    \end{remark}
    %
    With above assumption we can compute relative returns for each option given~$h_0$ as
    \eq{ret}
        \Delta H^\ell_\tau(h_0) := \frac1{ H^\ell(h_0) } \left(  H^\ell_\tau(h_0)  - H^\ell(h_0)  \right)
    \eqend
    We will use $\Delta H\in\R^{1+L}$ to refer to the vector of returns given~$\tau$.
    
	\begin{remark}
		For swaps the initial value $H_j^\ell$ is zero. In this case, we use either
		the value of the fixed leg (which is usually readily computed) or otherwise
		the ``notional''.	
		
		For futures the returns in each case refer to the accrued pnl from
		daily margin calls.
	\end{remark}


    \begin{assumptions}
    \hspace{1cm}
    	\begin{enumerate}
			\item Assume for now interest rates, 
			forward rates and dividends are zero.
    		\item $ \Delta H$ is a $C^2$ function of the state~$h_0$.
    		\item Assume that $\E^\P[| \Delta H(h_0)| ]<\infty$.
    	\end{enumerate}
    \end{assumptions}

    The main focus of this paper is the construction of risk-neutral margingale measures
    if they exist, and to ascertain which conditions are sufficient to do so in the setting
    introduced earlier. A well known key result due to~\cite{FFTOAP_DT} is 
	that existence of a martingale measure
	is equivalent to absence of arbitrage. We note here that absence of arbitrage
    is a \emph{local} condition for one return: an asset is a martingale if
    its conditional return is zero. 
    
	
	\begin{prop}[First Fundamental Theorem of Asset Pricing]	
		\hspace{1cm}\\
		Let $c\in\R_{\geq 0}^{1+:}$.
		Assume that there exists an arbitrage
		opportunity $g\in \R^{1+L}$ (conditional on~$h$)
		which means
		\eq{arb}
			g'\Delta H - |g'|c \geq 0 \ \ \ 
			\mbox{and} \ \ \ g'\Delta H - |g'|c \not=0 \ .
		\eqend 	
		Then no equivalent near-martingale or, if $c\equiv0$, martingale measure~$\Q$
		exists.
	\end{prop}



    Considering the Markov nature of our setup with~\eqref{h} with both~\eqref{D} and~\eqref{Delta}
    will often consider states and returns accross time jointly. This is justified by the following
    simplisitc proposition


    
    \begin{definition}
    	Let $u$ be a $C^1$ utility function which is strictly
    	concave, strictly increasing, is centered to $u(0)=0$
    	and $u'(0)=1$,\footnote{both can always be achieved
    	by translation for non-degenerate utilities} and
    	$u'(-\infty) = \infty$ and $u'(\infty) = 0$. 
    	
    	The Optimized Certainty Equivalent (OCE) of $u$ is
    	then defined as
    	\eqx
    		U(X) := \sup_{y\in\R} \E^\P\left[ u(y+X)-y \right]
    	\eqend
    	for any $L^1$ random variable $X$.
    \end{definition}
    
	\begin{example}
		The most important examples are:
		\begin{itemize}
			\item \textbf{Entropy}: the entropy for a risk-aversion $\lambda>0$ is given by
			\eqx
				u(x) := \frac{ 1-\exp(-\lambda,x) }\lambda
			\eqend
			for which the OCE has the explicit solution
			$$
				U(X) = -\frac1\lambda \log \E^\P\left[ \exp(-\lambda X)\right] \ .
			$$
			The entropy reduces to mean-variance for normal~$X$
			and in this and many other ways a natural mathematical choice
			for a utility. However, the utility of a short position
			in a Black \& Scholes equity is negative infinite which
			limits its theoretical appeal and in practise easily
			leads to exploding gradients in machine learning.
			
		    \item \textbf{Truncated Exponential}: we introduce
		    therefore the truncated $C^2$ exponential as
		    \eqx
		    	u(x) := \left\{
		    		\begin{array}{lll}
		    			\frac{1 - \exp(-\lambda\,x)}\lambda & x\geq 0 \\
		    			x - \frac\lambda2 (-x)^2 & x<0 	 
		    		\end{array}
		    	\right.
		    \eqend
		\end{itemize}
	\end{example}
	We note that our definition here is tighter than~\cite{DH} or~\cite{NEARMART} since 
	we require strict monotonicity and concavity, and because $u'(-\infty) = \infty$.
	Those conditions \emph{exclude} the CVaR OCE measure given by $u(x) = (1+\lambda) \min\{ 0, x \}$.
    
	We have shown before the following result:
	\begin{theorem}[Near-Martingale Measures~\cite{NEARMART}]\label{th:remove}\hspace{1cm}\\
		Let $c\in\R_{\geq 0}^{1+L}$.
		Let~$(a^*,y^*)$ be a finite optimal solution to the problem
		\eq{coeopt}
			\sup_{a\in\R^{1+L},y\in\R}:\ \E^\P\left[\ u\big( y + a' \Delta H
			- |a|'c
			\big)-y\ \right] 
		\eqend
        (conditional on~$h_0$)
		with
		$$
			(*) \ \ \ \E^\P\left[\ u\left( y^* + {a^*}' dH -
			|a^*|'c 
			\right)-y^*\ \right] < \infty \ .
		$$
		Then, then measure~$\Q$ 
		defined (conditional on~$h_0$) via the density
		\eqx
			\frac{d\Q}{d\P} := u'\big(\ y^* + {a^*}' \Delta H - |a^*|' c\ \big) 
		\eqend
		is a \emph{near-martingale measure} in the sense that
		$$
			\Big|\ \E^\Q\big[\ \Delta H\ \big]\ \Big| \leq c \ .
		$$
		Specifically, if $c\equiv 0$ then $\Q$ is a full martingale
		measure.
	\end{theorem}
	%
	This theorem yields a risk-neutral density if
	our optimal solution does not explode and satisfies~$(*)$. 	 (The actual statement in~\cite{NEARMART} is stronger
		and covers not only proportional transaction cost, but also
		more general convex cost.)
	The 
	open question so far was when that will happen. The
	following is the main result of this section (the existence of a~$\Q$
    with the specific construction above being the contribution here):
	
	\begin{prop}[Absence of Arbitrage for Near-Martingale Measures]
		\hspace{1cm}\\
		Let $c\in\R_{\geq 0}^{1+}$.
	Assume that the following integrability condition  
	\eq{intgrd}
		\E^\P\big[\, u\left(\,y + a\, \Delta H - |a|c\,\right)\, \big] < \infty 
	\eqend
	holds for all $y\in\R$ and $a\in\R^{1+L}$.

	Then a finite solution to~\eqref {coeopt} and therefore
	the [near-] martingale measure~$\Q$ exists if and only if
	there
	is no arbitrage opportunity which means for
	there is no $a\in\R^{1+L}$
	\eq{noarbcost}
			a'\Delta H - |a'|c \geq 0 \ \ \ \mbox{with} \ \ \ \P[ a'\Delta H - |a'|c ] > 0 \ .
	\eqend	
	In the case of zero transaction cost and absent any market frictions
	this condition becomes the classic
	\eq{noarb}
			a'\Delta H \geq 0 \ \ \ \mbox{with} \ \ \ \P[ a'\Delta H  ] > 0 \ .
	\eqend
	\end{prop}	
	For the utilities introduced above the integrability
	condition~\eqref {intgrd} is satisfied if:
	\begin{itemize}
		\item \textbf{Entropy}: all moment $\E^\P[ a'\Delta H ]<\infty$
		for $a\in\R^{1+L}$
		must exist.
		
		\item \textbf{Truncated Entropy}: the second moment
		must exist
		$\E^\P[ \Delta H^2 ]<\infty$.
	\end{itemize}

    \begin{prop}[SANOS]
        Assume that the market is given by SANOS and that $\tau\geq\max\{\tau^\ell\}$ (i.e.~the returns $dH$
        are computed to expiry of each option). Then~\eqref{noarb} is satisfied.
    \end{prop}

    \begin{proof}
        Absence of static arbitrage for each state~$h_0$ is exactly condition~\eqref{noarb}
        where for each expiry $\tau^\ell\leq \tau$ the spot takes as values the entire
        positive real line.
        
        Statistically, this statement is also true
        provided the spot at expiry covers the entire strike range of all options and beyond.        
    \end{proof}

	\noindent	
	We now translate the above into a setting suitable for machine
	learning:
	\begin{itemize}
		\item
			We normalize
			$$
				\Delta \hat H_j := \frac{ \Delta H_j }{ H_j } 
				\ \ \ \mbox{and} \ \ \ 
				D \hat H_j := \frac{ D H_j }{ H_j } 
			$$
			which represents in each case the return on a $1\$$
			investment.
			We simularly scale cost to
			$$
				\hat c_j := \frac{ c_j }{ H_j }\ .
			$$
			These definitions make sense if the trading instruments are ``funded''. 
			Swaps, on the other hand, have zero initial value as do futures. In such a case we use the value of the ``fixed leg`` 
			as denominator (swaps) or its quoted spot value (futures).
		\item
			Accordingly, an action $a$ with payoff 
			$$
				a' d\hat H_j - |a'|\hat c_j 
			$$
			now represents heuristically ``capital invested''.
			We parameterize an action therefore as
			$$
				a_j^\ell := a^{\max}_j \alpha^\ell_j
			$$
			with
			\begin{itemize}
				\item maximum cash invested $a^{\max}\in\R$ and
				\item its allocation $\alpha_j \in (-1,+1)^{1+L}$.t
			\end{itemize}
	\end{itemize}



	\pagebreak    
    
    Let
    $$
    		\sup_{y,a} : - \frac1\lambda\log  \E_b\left[ \exp\left( - \lambda \E_\omega\left[ D_b: u(y + a\,dX) - y \right] \right) \right]
	$$
    $$
    		\sup_{y,a} : - \E_b\left[ \lambda \exp\left( - \lambda \E_\omega\left[ D_b: u(y + a\,dX) - y\right] \right) \right]
	$$
	$$
		\partial_a = \E_b\left[ \exp\left( - \E_\omega\left[ D_b: u(y + a\,dX) - y\right] \right) \E_\omega\left[ D_b u'(y + a\,dX)\,dX \right] \right]
	$$
	$$
		\partial_a = \E_{b\times\omega} \left[ \exp\left( - \E_\omega\left[ D_b: u(y + a\,dX) - y\right] \right)\, D_b\,u'(y + a\,dX)\,dX  \right]
	$$
	$$
		\partial_a = \E_{\omega} \left[ \E_b\left\{ D_b 
					\exp\left( - \E_\omega\left[D_b: u(y + a\,dX) - y\right] \right)\right\}\, u'(y + a\,dX)\,dX  \right]
	$$
    
    


    


	\pagebreak    
    
    $$

    
    % ------------------------------------------------------------------
    % ------------------------------------------------------------------

   \addcontentsline{toc}{section}{Bibliography}

   \bibliography{references}
   \bibliographystyle{alpha}
  


\end{document}/